% Original coefficient reduction. All maps are complex-linear until a metric is supplied.
\section{An exact conductor frame and a fixed invariant-generator recursion}
\label{sec:OCF}

Retain the original simple-quartet data, $m=1$, $k=4l+1$,
$l\geq1$, the original root order $(++),(+-),(-+),(--)$, and
the complete original period and amplitude matrix of OQC1:
\[
                  \mathsf H=F^{-1}\Pi U\mathcal E_{\rm prim}.
\]
Thus $\det\mathsf H\ne0$ and
$Q(x)=Q_0(\mathsf H^{-1}x)$, where
$Q_0(t)=t_{++}t_{--}-t_{+-}t_{-+}$.
Use $g f(x)=f(D^{-1}x)$,
$D=\operatorname{diag}(\zeta,\zeta^2,\zeta^3,\zeta^4)$,
$\zeta=e^{2\pi i/5}$, and the original invariant algebra
$\mathscr I=\mathbb C[x_1,x_2,x_3,x_4]^{\langle g\rangle}$.
This calculation applies on the actual five-orbit period locus,
which includes every original logarithm branch with $|u|\geq R_*$
by PQO and RPC. On that locus the five prime polynomials
$Q_a=g^aQ$, $0\leq a<5$, are pairwise nonassociate. Keep their
actual products
\[
 \mathcal A=\prod_{a=1}^4Q_a,\qquad
 \mathcal H_Q=Q\mathcal A=\prod_{a=0}^4Q_a,
 \qquad \deg\mathcal A=8,\quad\deg\mathcal H_Q=10.
 \tag{OCF1}
\]
No coefficient in these polynomials is replaced. The exact orbit
decision and the period domain are those of the cited original
calculations. The construction below treats every possible nonzero
leading coefficient of their actual restriction.

\subsection{The original Segre coefficient isomorphism}

For $d\geq0$ let $\mathcal B_d$ be the space of biforms of
bidegree $(d,d)$ in $(z_0,z_1)$ and $(w_0,w_1)$. Its ordered basis is
\[
 e_{ab}^{(d)}=z_0^a z_1^{d-a}w_0^b w_1^{d-b},
 \qquad 0\leq a,b\leq d.
 \tag{OCF2}
\]
Pairs are ordered lexicographically, first by $a$, then by $b$.
The order is used only for the exact coefficient elimination.
Set $\mathcal B_d=0$ for $d<0$. Multiplication satisfies
$e_{ab}^{(d)}e_{ij}^{(e)}=e_{a+i,b+j}^{(d+e)}$.
The original substitution is
\[
 \begin{gathered}
 t=(z_0w_0,z_0w_1,z_1w_0,z_1w_1)^T,\qquad x=\mathsf Ht,\\
 \mathscr P_d=\mathbb C[x_1,x_2,x_3,x_4]_d,\qquad
 \psi_d:\mathscr P_d\longrightarrow\mathcal B_d,\quad
 f\longmapsto f(\mathsf Ht).
 \end{gathered}\tag{OCF3}
\]
The induced map
$R_d=\mathscr P_d/(Q\mathscr P_{d-2})\longrightarrow\mathcal B_d$
is an isomorphism. Here is a coefficient proof with the given
orientation. For each $(a,b)$, choose an integer
$h$ between $\max(0,a+b-d)$ and $\min(a,b)$. The monomial
$t_{++}^{h}t_{+-}^{a-h}t_{-+}^{b-h}t_{--}^{d-a-b+h}$ maps
to $e_{ab}^{(d)}$, proving surjectivity. Two monomials with
the same pair differ in their exponent vectors by an integer
multiple of $(1,-1,-1,1)$. After removing their common monomial,
their difference is divisible by $t_{++}t_{--}-t_{+-}t_{-+}$,
using $A^n-B^n=(A-B)\sum_{j=0}^{n-1}A^{n-1-j}B^j$.
A polynomial in the kernel has coefficient sum zero within every
pair, so it is a sum of such differences. The reverse inclusion
in the kernel follows by direct substitution. The invertible
substitution $x=\mathsf Ht$ then gives exactly the asserted kernel.
Thus these maps preserve multiplication as well as every grading.

Write the full conductor restriction as
\[
 A=\psi_8(\mathcal A)
     =\sum_{r=0}^8\sum_{s=0}^8 a_{rs}e_{rs}^{(8)}\ne0.
 \tag{OCF4}
\]
To verify its nonvanishing, $Q$ is a rank-four irreducible quadric,
hence prime in the polynomial ring, and none of $Q_1,\ldots,Q_4$
is its associate. Therefore $Q$ does not divide their product.
The kernel calculation OCF3 proves OCF4. The ring
$\mathcal B=\bigoplus_{d\geq0}\mathcal B_d$, with the indicated
diagonal grading, embeds in the four-variable polynomial ring;
it is an integral domain. Multiplication by $A$ is consequently
injective on every $\mathcal B_d$.

\subsection{Every leading-coefficient case gives an exact strip frame}

Let $(r_*,s_*)$ be the lexicographically largest pair for which
$a_{r_*s_*}\ne0$, and put $a_*=a_{r_*s_*}$. This is an
actual nonzero coefficient of OCF4. For $d\geq8$ define
\[
 \begin{split}
 P_d&=\{(r_*+i,s_*+j):0\leq i,j\leq d-8\},\\
 E_d^{\rm strip}&=\{0,\ldots,d\}^2\setminus P_d,\qquad
 W_d=\mathbb C^{E_d^{\rm strip}}.
 \end{split}\tag{OCF5}
\]
For $0\leq d<8$ put $P_d=\varnothing$,
$E_d^{\rm strip}=\{0,\ldots,d\}^2$ and use the same definition
of $W_d$; for $d<0$ put $W_d=0$. Let
$S_d:W_d\hookrightarrow\mathbb C^{(d+1)^2}$ insert zeros in
the positions $P_d$. All coefficient columns below use OCF2.

The following finite elimination defines
$N_d:\mathbb C^{(d+1)^2}\to W_d$ and, for $d\geq8$,
$V_d:\mathbb C^{(d+1)^2}\to\mathbb C^{(d-7)^2}$.
Starting with the complete input biform $f$, traverse all
$(i,j)\in\{0,\ldots,d-8\}^2$ in decreasing lexicographic order.
At the step $(i,j)$, let $c$ be the current coefficient at
$(r_*+i,s_*+j)$, record $b_{ij}=c/a_*$, and replace the current
biform by
\[
       f_{\rm current}-b_{ij}A e_{ij}^{(d-8)}.
 \tag{OCF6}
\]
At completion take the coefficients on $E_d^{\rm strip}$ as
$N_df$ and take $(b_{ij})$ as $V_df$. If $d<8$, $N_d$ is the
identity and $V_d$ is the zero map into the zero space.
The exact multiplication matrix $M_{A,d}$ gives
\[
 S_dN_d+M_{A,d}V_d=I,\qquad
 N_dS_d=I_{W_d},\qquad
 \ker N_d=A\mathcal B_{d-8}.
 \tag{OCF7}
\]
To prove these identities, the largest term of
$A e_{ij}^{(d-8)}$ is $a_*e_{r_*+i,s_*+j}^{(d)}$;
every other term is strictly smaller in the same lexicographic
order. Therefore each step zeros its designated position without
altering an earlier zero. Summing all subtractions gives the
first identity. Input supported on the strip produces zero
coefficients at every elimination step, giving the second.
Finally, the leading term of $Ab$, for a nonzero homogeneous
biform $b$, is the product of the two leading terms and belongs
to $P_d$. Indeed every other product is smaller in the additive
lexicographic order. Thus
$A\mathcal B_{d-8}\cap S_dW_d=0$. The first identity now proves
the kernel formula, the uniqueness of the decomposition, and all
three identities in OCF7. All operations are linear in the input;
their coefficients are explicit rational expressions in the
original $a_{rs}$ with powers of the displayed $a_*$ as denominators.
The polynomial $A$ itself remains the full polynomial OCF4.

In particular the quotient map induced by $N_d$ is the exact
isomorphism
\[
 \mathcal B_d/(A\mathcal B_{d-8})\xrightarrow{\ \sim\ }W_d,
 \qquad
 n_d:=\dim W_d=
 \begin{cases}(d+1)^2,&0\leq d<8,\\16d-48,&d\geq8.
 \end{cases}
 \tag{OCF8}
\]
The second formula is
$(d+1)^2-(d-7)^2=16d-48$. The possible choices
$(r_*,s_*)\in\{0,\ldots,8\}^2$ form an exhaustive finite
partition: the chart for a pair specifies that its coefficient
is nonzero and every larger coefficient is zero. Hence all actual
coefficient cases have been included. At the first original
degree $k=5$ this construction retains all $36$ coordinates;
at every original $k\geq9$ it retains $16k-48$ coordinates.

\subsection{Twenty-five fixed invariant monomials generate all degrees}

Define the following fixed finite sets of exponent vectors:
\[
 G_e=\{\nu\in\mathbb N^4:|\nu|=e,
                   \nu_1+2\nu_2+3\nu_3+4\nu_4\equiv0\pmod5\},
 \quad 1\leq e\leq5,\qquad G=\bigcup_{e=1}^5G_e.
 \tag{OCF9}
\]
Here $G_1$ is empty. The complete remaining lists, in the original
$x_1,x_2,x_3,x_4$ order, are
\[
\begin{array}{c|l}
e&G_e\\\hline
2&(0,1,1,0),(1,0,0,1)\\
3&(0,0,2,1),(0,1,0,2),(1,2,0,0),(2,0,1,0)\\
4&(0,0,1,3),(0,2,2,0),(0,3,0,1),(1,0,3,0),\\
 & (1,1,1,1),(2,0,0,2),(3,1,0,0)\\
5&(0,0,0,5),(0,0,5,0),(0,1,3,1),(0,2,1,2),\\
 & (0,5,0,0),(1,0,2,2),(1,1,0,3),(1,3,1,0),\\
 & (2,1,2,0),(2,2,0,1),(3,0,1,1),(5,0,0,0).
\end{array}\tag{OCF10}
\]
Thus $(|G_1|,\ldots,|G_5|)=(0,2,4,7,12)$ and $|G|=25$.
These lists can also be obtained directly by choosing
$0\leq\nu_1,\nu_2,\nu_3\leq e$ with their sum at most $e$,
putting $\nu_4=e-\nu_1-\nu_2-\nu_3$, and applying the single
displayed congruence, which verifies completeness of each list.

Every positive-degree invariant monomial is a product of members
$x^\nu$, $\nu\in G$. To prove this, list the weights of its
individual factors in any order. If its degree is at most five,
the entire list gives a member of $G$. If its degree exceeds five,
consider the six partial sums of the first five weights, including
the empty partial sum, in $\mathbb Z/5\mathbb Z$. Two coincide;
the nonempty block between them has length at most five and total
weight zero. Its factors give a divisor $x^\nu$ with $\nu\in G$.
The complementary monomial again has total weight zero and smaller
degree. Induction completes the factorization. Because the action
is diagonal, a polynomial is invariant exactly when all its
nonzero monomials have weight zero: comparison of each monomial
coefficient in $gf=f$ proves this statement. Thus
\[
                \mathscr I=\mathbb C[x^\nu:\nu\in G].
 \tag{OCF11}
\]
All intermediate degrees here are degrees in the original
coefficient algebra. This construction introduces no new source
distribution and leaves the original source orders $1$ and $k$ intact.

For each $\nu\in G_e$ the complete generator biform is
\[
 f_\nu=\psi_e(x^\nu)
  =\prod_{j=1}^4
   (H_{j1}z_0w_0+H_{j2}z_0w_1+
                  H_{j3}z_1w_0+H_{j4}z_1w_1)^{\nu_j}
  =\sum_{a,b=0}^e c_{\nu,ab}e_{ab}^{(e)}.
 \tag{OCF12}
\]
For an entirely explicit coefficient formula, sum over all
nonnegative $4$ by $4$ integer arrays $(n_{jt})$ satisfying
$\sum_t n_{jt}=\nu_j$,
$\sum_j(n_{j1}+n_{j2})=a$ and
$\sum_j(n_{j1}+n_{j3})=b$. Then
\[
 c_{\nu,ab}=
 \sum_{(n_{jt})}
     \prod_{j=1}^4\left(
        \frac{\nu_j!}{\prod_{t=1}^4n_{jt}!}
        \prod_{t=1}^4H_{jt}^{n_{jt}}\right).
 \tag{OCF13}
\]
The multinomial theorem proves the formula, including every
original period and amplitude entry. Every factor has degree at
most five and at most $36$ possible biform coefficients.

\subsection{Exact multiplication on the strip and the invariant recursion}

Let $\mathcal L=\mathcal B/(A\mathcal B)$, a graded algebra.
For $\nu\in G_e$ and $d\geq e$ define
\[
 m_{\nu,d}:W_{d-e}\longrightarrow W_d,
 \qquad m_{\nu,d}=N_d M_{f_\nu}S_{d-e}.
 \tag{OCF14}
\]
Here $M_{f_\nu}$ is the exact coefficient convolution:
\[
       (M_{f_\nu}S_{d-e}v)_{ab}
          =\sum_{i,j}c_{\nu,a-i,b-j}(S_{d-e}v)_{ij},
\]
with out-of-range indices giving zero.
Therefore every entry of OCF14 is supplied
by the finite formulas OCF6 and OCF13.
For a representative $h\in\mathcal B_{d-e}$, OCF7 writes
$h-S_{d-e}N_{d-e}h=A b$. Multiplication by $f_\nu$ preserves
this ideal. Applying $N_d$ yields the exact identity
\[
       N_d(f_\nu h)=m_{\nu,d}N_{d-e}h.
 \tag{OCF15}
\]
This proves that OCF14 is multiplication by the specified element
of $\mathcal L$, including its independence of representatives.

The original invariant image is
$J_d=\psi_d(\mathscr I_d)\subseteq\mathcal B_d$.
Set $\overline J_d=N_dJ_d\subseteq W_d$, and set
$\overline J_d=0$ for $d<0$. The exact recursion is
\[
 \overline J_0=\mathbb C\cdot1,\qquad
 \overline J_d=
   \sum_{\substack{\nu\in G_e\\2\leq e\leq5,\ e\leq d}}
                 m_{\nu,d}\overline J_{d-e}
                 \quad(d\geq1).
 \tag{OCF16}
\]
Indeed OCF11 factors every invariant monomial of degree $d>0$
as $x^\nu$ times an invariant monomial of degree $d-e$.
Applying $\psi$, then OCF15, gives the inclusion from left to
right. Conversely every summand is the image of a product of two
invariant polynomials, so it belongs to the left side. Linearity
proves equality for the complete invariant space. In particular
the recursion gives $\overline J_1=0$ and includes all small degrees.

The conductor ideal is contained in this invariant image with
its original trace coefficient:
\[
 A\mathcal B_{d-8}\subseteq J_d,\qquad
 [5\mathsf R(\mathcal A h)]_{(Q)}=[\mathcal A h]_{(Q)},
 \quad\mathsf R=\frac15\sum_{a=0}^4g^a.
 \tag{OCF17}
\]
To prove the identity, $g^a\mathcal A$ is the product of all
five orbit equations except $Q_a$; for $a\ne0$ it contains $Q$.
Thus all four corresponding terms vanish modulo $Q$, and the
$a=0$ term is the displayed product. Every biform in
$\mathcal B_{d-8}$ has a polynomial lift by OCF3, so the identity
proves the inclusion. It also shows exactly why the quotient
used in OCF16 retains the original invariant image.

\subsection{Exact ranks, a finite pivot atlas, and the original quotient map}

For $d\geq0$ put
\[
 d_0(d)=\frac{\binom{d+3}{3}+4\epsilon_d}{5},\qquad
 \epsilon_d=\mathbf1_{d\equiv0\ (5)}-\mathbf1_{d\equiv1\ (5)},
 \qquad d_0(d)=0\quad(d<0).
 \tag{OCF18}
\]
This is the number of degree-$d$ weight-zero monomials: the
identity character has generating function $(1-z)^{-4}$, each
of the four nonidentity characters has generating function
$\prod_{j=1}^4(1-z\zeta^j)^{-1}=(1-z)/(1-z^5)$, and averaging
their coefficients gives OCF18.
Every invariant polynomial divisible by $Q$ is divisible by all
five $Q_a$, hence by $\mathcal H_Q$, since they are pairwise
nonassociate primes. Its quotient by this product is invariant:
apply $g$ and cancel the nonzero invariant product in the integral
domain. Conversely each $\mathcal H_Q$ times an invariant
polynomial is invariant and vanishes modulo $Q$. Consequently
$\ker(\mathscr I_d\to J_d)=\mathcal H_Q\mathscr I_{d-10}$,
and injection of multiplication by the nonzero product gives
$\dim J_d=d_0(d)-d_0(d-10)$.
Injection of multiplication by $A$, OCF7 and OCF17 then prove
\[
 j_d:=\dim\overline J_d
     =d_0(d)-d_0(d-10)-\max(d-7,0)^2.
 \tag{OCF19}
\]
In particular $j_0=1$, $j_1=0$, $j_5=12$, and
$j_d=8d-32$ for $d\geq8$. For the last formula at $d\geq10$,
$\epsilon_d=\epsilon_{d-10}$ and direct expansion gives
$\binom{d+3}{3}-\binom{d-7}{3}=5(d^2-6d+17)$;
subtracting $(d-7)^2$ gives $8d-32$. At $d=8,9$ the ranks
before this subtraction are $33,44$, respectively, which gives
$32,40$ and verifies both remaining cases.

Here is a complete finite construction of all the columns used
by OCF16. Start with $Z_0=(1)$ and the invariant polynomial
certificate $F_{0,1}=1$. Suppose all earlier matrices $Z_b$ have
independent columns spanning $\overline J_b$, and their columns
are $N_b\psi_b(F_{b,h})$ for invariant polynomials $F_{b,h}$.
Form the candidate matrix
\[
 C_d=\bigl[\,m_{\nu,d}Z_{d-e}\,\bigr]_{
                    \nu\in G_e,\ 2\leq e\leq5,\ e\leq d}.
 \tag{OCF20}
\]
Each candidate column has the explicit invariant certificate
$x^\nu F_{d-e,h}$. OCF16 proves that $C_d$ has column image
$\overline J_d$ and rank $j_d$. Select $j_d$ columns and $j_d$
rows giving a nonzero determinant. Let $Z_d$ be those columns,
let $I_d$ be those row positions, and let $T_d$ be the complementary
row positions in $E_d^{\rm strip}$. Retain the corresponding
invariant certificates. This completes the induction. When
$j_d=0$, use the empty matrix with zero columns, empty $I_d$,
and $T_d=E_d^{\rm strip}$.

Every selected determinant and every matrix entry is an explicit
expression obtained from the original entries of $\mathsf H$ and
the full $a_{rs}$. To cover all cases, order the finite choices
of column and row subsets and select the first nonzero determinant;
the corresponding chart asserts its nonvanishing and the vanishing
of all preceding candidates. A nonzero determinant exists because
the rank was proved in OCF19. At rank zero no determinant is needed.
Together with the exhaustive leading-coefficient charts after
OCF8, this gives a finite atlas for each specified $k$. It requires
no generic coefficient premise and supplies no numerical zero test
for an unevaluated analytic coefficient: each chart is specified
by its exact minors.

Define $\rho_d:W_d\to\mathbb C^{T_d}$ and the full original
coefficient map $\pi_d$ by
\[
 \rho_d=\operatorname{row}_{T_d}
            -(Z_d)_{T_d}(Z_d)_{I_d}^{-1}
                                  \operatorname{row}_{I_d},
 \qquad \pi_d=\rho_dN_d.
 \tag{OCF21}
\]
At rank zero the second term is the zero map. If $E_{T_d}$
inserts a column into its $T_d$ positions in $W_d$, then
\[
 \begin{split}
 w&=Z_d(Z_d)_{I_d}^{-1}w_{I_d}+E_{T_d}\rho_dw,\\
 \rho_dE_{T_d}&=I,\qquad\ker\rho_d=\operatorname{im}Z_d
                                  =\overline J_d,\\
 \ker\pi_d&=J_d,\qquad
       \pi_d S_d E_{T_d}=I_{\mathbb C^{T_d}}.
 \end{split}\tag{OCF22}
\]
For the first line, the subtraction on its right leaves zero
$I_d$ rows and exactly $\rho_dw$ on the complementary rows.
This proves the first three claims. If $\pi_df=0$, then
$N_df\in\overline J_d$; choose $j\in J_d$ with $N_dj=N_df$.
Now $f-j\in\ker N_d=A\mathcal B_{d-8}\subseteq J_d$ by
OCF17, so $f\in J_d$. The reverse inclusion follows immediately
from $N_dJ_d=\overline J_d$. The last equality uses $N_dS_d=I$.
This proves every map and kernel in OCF22.

At the original degrees, these are the exact dimensions
\[
 \begin{array}{c|c|c|c}
 k&n_k&j_k&|T_k|=n_k-j_k\\\hline
 5&36&12&24\\
 k\geq9,\ k=4l+1&16k-48&8k-32&8k-16.
 \end{array}\tag{OCF23}
\]
The original map $\Theta_k:R_k\xrightarrow{\sim}E_k^\vee$
in OQC7 sends the coefficient basis OCF2 to the original dual
primary basis in the same order. Indeed its evaluation on
$e^{wY_k}[1]$ is
$\exp(((2b-k)\gamma-i(2a-k)\delta)w)$. The corresponding
eigenvalue of $A_k$ is
$\beta_{ab}=k/2+(2a-k)\delta+i(2b-k)\gamma$; the eigenvalue
of $Y_k=(A_k-kI/2)/i$ is
$\omega_{ab}=(\beta_{ab}-k/2)/i=(2b-k)\gamma-i(2a-k)\delta$.
These eigenvalues are distinct because
$\delta>0$ and $\gamma>0$. Their exponentials are independent
by the Vandermonde determinant, so the evaluation identifies the
covectors exactly. By the original observation identity OQC8,
$\Theta_kJ_k=\operatorname{im}\Lambda_k^\vee$. Therefore OCF22
supplies the specified isomorphism
\[
 E_k^\vee/\operatorname{im}\Lambda_k^\vee
       \xrightarrow{\ \pi_k\ }\mathbb C^{T_k}
       \cong (\ker\Lambda_k)^\vee.
 \tag{OCF24}
\]
The displayed $\pi_k$ uses exactly the original primary coefficient
coordinates under $\Theta_k$. Every dualization in this statement
is complex-linear. No conjugation is inserted in this coefficient
map; the original metric maps can subsequently be composed with it.

Finally, this construction is finite with explicitly bounded
compressed matrices. OCF20 has $n_d$ rows and
$\sum_{e=2}^5 |G_e|j_{d-e}$ columns, with negative subscripts zero.
Both dimensions grow at most linearly in $d$: for example
$n_d\leq16d+1$ and the column number is at most $25(16d+1)$.
For $d\geq13$ the exact column number is $200d-1632$, obtained
from $\sum_e|G_e|=25$, $\sum_e e|G_e|=104$, and
$j_{d-e}=8(d-e)-32$. Only the preceding five graded bases are
needed to form the next matrix. Each column is obtained from
the at-most-$36$-coefficient formula OCF12, one exact convolution,
and OCF6. A full original biform coefficient column, if used as
temporary storage, has $(d+1)^2$ entries; every invariant relation
matrix in OCF20 has the stated linear row and column counts.
The invariant certificates can be retained as their products and
parent-column references, so their expanded high-degree monomials
need never be enumerated. The proof of their completeness is
OCF11--16, and the exact original quotient and its right inverse
are OCF21--24.

\subsection{The same recursion with fourteen proved generators}

The complete fixed set OCF9 can be reduced to the following
fourteen specific members while preserving every map and image
just proved:
\[
\begin{array}{c|l}
e&\mathcal G_e\\\hline
2&x_1x_4,\ x_2x_3\\
3&x_1^2x_3,\ x_1x_2^2,\ x_2x_4^2,\ x_3^2x_4\\
4&x_1^3x_2,\ x_1x_3^3,\ x_2^3x_4,\ x_3x_4^3\\
5&x_1^5,\ x_2^5,\ x_3^5,\ x_4^5 .
\end{array}\tag{OCF25}
\]
Here is a complete proof. Call a positive-degree invariant
monomial indecomposable when it is not the product of two
positive-degree invariant monomials. The divisor proof preceding
OCF11 shows that each indecomposable has degree at most five.
If it contains both $x_1,x_4$, it is their product, since otherwise
division by this invariant quadratic leaves a nonconstant
invariant monomial. The same assertion holds for $x_2,x_3$.
After these cases are removed, at most two variables occur:
every three indices contain one of the two opposite pairs.
One variable gives precisely its fifth power. The remaining
two-element supports are $\{1,2\},\{1,3\},\{2,4\},\{3,4\}$.
For positive exponents $(a,b)$ with $a+b\leq5$, the equation
$a+2b\equiv0\pmod5$ has exactly $(a,b)=(3,1),(1,2)$, and
$a+3b\equiv0\pmod5$ has exactly $(a,b)=(2,1),(1,3)$.
For $\{2,4\}$ divide the weights by $2$ in the field
$\mathbb Z/5\mathbb Z$, giving the first list; for $\{3,4\}$
divide by $3$, giving the second. These are exactly the eight
mixed monomials of degrees three and four in OCF25. Every
listed monomial has no proper invariant divisor, as is also
seen from these lists and the excluded opposite pairs.
This proves the classification. Factoring any invariant
monomial repeatedly into indecomposables terminates because
positive total degree strictly decreases. Hence OCF25 generates
the entire invariant algebra.

Use $\mathcal G=\bigcup_{e=2}^5\mathcal G_e$ in place of
$\{x^\nu:\nu\in G\}$ in OCF12--16 and OCF20. The same
factorization argument proves the same exact invariant image;
all coefficient formulas, invariant certificates, pivot charts,
quotient maps and ranks remain valid. The new recursion has
$2,4,4,4$ generators in degrees $2,3,4,5$, respectively. Thus
for $d\geq13$ its precise matrix dimensions and rank are
\[
 C_d^{(14)}\in
 \operatorname{Mat}_{(16d-48)\,\times\,(112d-864)}(\mathbb C),
 \qquad \operatorname{rank}C_d^{(14)}=8d-32.
 \tag{OCF26}
\]
Indeed the number of generators is $14$ and the sum of their
degrees is $2\cdot2+4\cdot3+4\cdot4+4\cdot5=52$; substituting
the already proved $j_{d-e}=8(d-e)-32$ gives
$14(8d-32)-8\cdot52=112d-864$. At smaller degrees the exact
number of columns is the same sum of the appropriate $j_{d-e}$
from OCF19, with negative subscripts zero. All coefficients
in this reduced recursion are still the original coefficients
of OCF12--14.
